{{دیوان
| عنوان = بخش ۲ - تمهید
| ماخذ = گنجور
| ماخذ_لنک = https://ganjoor.net/iqbal/asrar-khodi/sh2
| مقابلہ = اقبال سائبر لائبریری (اقبال اکادمی پاکستان)
| مقابلہ_لنک = https://www.iqbalcyberlibrary.net
}}
<poem>
نیست در خشک و تر بیشهٔ من کوتاهی
چوب هر نخل که منبر نشود دار کنم

نظیری نیشابوری

راه شب چون مهر عالم‌تاب زد
گریهٔ من بر رخ گل، آب زد

اشک من از چشم نرگس خواب شُست
سبزه از هنگامه‌ام بیدار رُست

باغبان زور کلامم آزمود
مصرعی کارید و شمشیری درود

در چمن جز دانهٔ اشکم نکِشت
تار افغانم به پود باغ رِشت

ذره‌ام مهر منیر آنِ من است
صد سحر اندر گریبان من است

خاک من روشن تر از جام جم است
محرم از نازاد‌های عالم است

فکرم آن آهو سرِ فتراک بست 
کو هنوز از نیستی بیرون نجست

سبزه ناروئیده زیب گلشنم
گل به شاخ اندر نهان در دامنم

محفل رامشگری برهم زدم
زخمه بر تار رگ عالم زدم

بس که عود فطرتم نادر نواست
هم‌نشین از نغمه‌ام نا آشناست

در جهان خورشید نوزائیده‌ام
رسم و آئین فلک نادیده‌ام

رَم ندیده انجم از تابم هنوز
هست ناآشفته سیمابم هنوز

بحر از رقص ضیایم بی‌نصیب
کوه از رنگ حنایم بی‌نصیب

خوگر من نیست چشم هست و بود
لرزه بر تن خیزم از بیم نمود

بامم از خاور رسید و شب شکست
شبنم نو بر گُل عالم نشست

انتظار صبح‌خیزان می کشم
ای خوشا زرتشتیان آتشم

نغمه‌ام از زخمه بی‌پرواستم
من نوای شاعر فرداستم

عصر من دانندهٔ اسرار نیست
یوسف من بهر این بازار نیست

ناامید استم ز یاران قدیم
طور من سوزد که می‌آید کلیم

قلزم یاران چو شبنم بی‌خروش
شبنم من مثل یم طوفان‌به‌دوش

نغمهٔ من از جهانِ دیگر است
این جرس را کاروان دیگر است

ای بسا شاعر که بعد از مرگ زاد
چشم خود بر بست و چشم ما گشاد

رخت باز از نیستی بیرون کشید
چون گل از خاک مزار خود دمید

کاروان ها گرچه زین صحرا گذشت
مثل گام ناقه کم‌غوغا گذشت

عاشقم. فریاد، ایمان من است
شور حشر از پیش‌خیزان من است

نغمه‌ام ز اندازهٔ تار است بیش
من نترسم از شکست عود خویش

قطره از سیلاب من بیگانه به
قلزم از آشوب او دیوانه به

در نمی‌گنجد به جو عمان من
بحرها باید پی طوفان من

غنچه کز بالیدگی گلشن نشد
در خور ابر بهار من نشد

برق‌ها خوابیده در جان من است
کوه و صحرا باب جولان من است

پنجه کن با بحرم ار صحراستی
برق من در گیر اگر سیناستی 

چشمهٔ حیوان براتم کرده‌اند
محرم راز حیاتم کرده‌اند

ذره از سوز نوایم زنده گشت
پر گشود و کرمک تابنده گشت

هیچکس، رازی که من گویم نگفت
همچو فکر من دُر معنی نسفت

سِرّ عیش جاودان خواهی بیا
هم زمین هم آسمان خواهی، بیا

پیر گردون با من این اسرار گفت
از ندیمان رازها نتوان نهفت

ساقیا برخیز و می در جام کن
محو از دل کاوش ایام کن

شعلهٔ آبی که اصلش زمزم است
گر گدا باشد پرستارش جم است

مِی کند اندیشه را هشیارتر
دیدهٔ بیدار را بیدارتر

اعتبارِ کوه بخشد کاه را
قوّت شیران دهد روباه را

خاک را اوج ثریا می‌دهد
قطره را پهنای دریا می‌دهد

خامشی را شورش محشر کند
پای کبک از خونِ باز احمر کند

خیز و در جامم شراب ناب ریز
بر شب اندیشه‌ام مهتاب ریز

تا سوی منزل کشم آواره را
ذوق بیتابی دهم نظّاره را

گرم‌رو از جستجوی نو شوم
روشناس آرزوی نو شوم

چشم اهل ذوق را مردم شوم
چون صدا در گوش عالم گم شوم

قیمت جنس سخن بالا کنم
آب چشم خویش در کالا کنم

باز بر خوانم ز فیض پیر روم
دفتر سربسته اسرار علوم

جان او از شعله‌ها سرمایه‌دار
من فروغ یک نفس مثل شرار

شمع سوزان تاخت بر پروانه‌ام
باده شبخون ریخت بر پیمانه‌ام

پیر رومی خاک را اکسیر کرد
از غبارم جلوه‌ها تعمیر کرد

ذره از خاک بیابان رخت بست
تا شعاع آفتاب آرد به دست

موجم و در بحر او منزل کنم
تا دُر تابنده‌ای حاصل کنم

من که مستی‌ها ز صهبایش کنم
زندگانی از نَفَس‌هایش کنم

شب دل من مایل فریاد بود
خامشی از «یا ربم» آباد بود

شکوهٔ شوب غم دوران بدم
از تهی‌پیمانگی نالان بدم

این قدر نظّاره‌ام بیتاب شد
بال و پر بشکست و آخر خواب شد

روی خود بنمود پیر حق‌سرشت
کو به حرف پهلوی قرآن نوشت

گفت «ای دیوانهٔ ارباب عشق
جرعه‌ای گیر از شراب ناب عشق

بر جگر هنگامهٔ محشر بزن
شیشه بر سر، دیده بر نشتر بزن

خنده را سرمایهٔ صد ناله ساز
اشک خونین را جگر پرکاله ساز

تا به کِی چون غنچه می‌باشی خموش
نکهت خود را چو گل ارزان فروش

در گره هنگامه داری چون سپند
محمل خود بر سر آتش ببند

چون جرس آخر ز هر جزوِ بدن
نالهٔ خاموش را بیرون فکن

آتش استی بزم عالم بر فروز
دیگران را هم ز سوز خود بسوز

فاش گو اسرار پیر مِی‌ فروش
موج می‌شو کسوت مینا بپوش

سنگ شو آئینهٔ اندیشه را
بر سر بازار بشکن شیشه را

از نیستان همچو نی پیغام ده
قیس را از قوم «حی» پیغام ده

ناله را انداز نو ایجاد کن
بزم را از های و هو آباد کن

خیز و جانِ نو بده هر زنده را
از «قم» خود زنده‌تر کن زنده را

خیز و پا بر جادهٔ دیگر بنه
جوش سودای کهن از سر بنه

آشنای لذت گفتار شو
ای دِرای کاروان بیدار شو»

زین سخن آتش به پیراهن شدم
مثل نِی هنگامهٔ بستن شدم

چون نوا از تار خود برخاستم
جنتی از بهر گوش آراستم

بر گرفتم پرده از راز خودی
وا نمودم سِرّ اعجاز خودی

بود نقش هستیَم انگاره‌ای
ناقبولی، ناکسی، ناکاره‌ای

عشق سوهان زد مرا، آدم شدم
عالِم کیف و کم عالم شدم

حرکت اعصاب گردون دیده‌ام
در رگ مه گردش خون دیده‌ام 

بهر انسان چشم من شب‌ها گریست
تا دریدم پردهٔ اسرارِ زیست

از درون کارگاه ممکنات
بر کشیدم سِر تقویم حیات

من که این شب را چو مه آراستم
گرد پای ملت بیضاستم

ملتی در باغ و راغ آوازه‌اش
آتش دل‌ها سرود تازه‌اش

ذره کِشت و آفتاب انبار کرد
خرمن از صد رومی و عطار کرد

آه گرمم، رخت بر گردون کشم
گرچه دودم از تبار آتشم

خامه‌ام از همت فکر بلند
راز این نُه پرده در صحرا فکند

قطره تا همپایهٔ دریا شود
ذره از بالیدگی صحرا شود

شاعری زین مثنوی مقصود نیست
بت پرستی، بتگری مقصود نیست

هندیم از پارسی بیگانه‌ام
ماه نو باشم تهی پیمانه‌ام

حُسن اندازِ بیان از من مجو
خوانسار و اصفهان از من مجو

گرچه هندی در عذوبت شکّر است
طرز گفتار دری شیرین‌تر است

فکر من از جلوه‌اش مسحور گشت
خامهٔ من شاخ نخل طور گشت

پارسی از رفعت اندیشه‌ام 
در خورَد با فطرت اندیشه‌ام

خرده بر مینا مگیر ای هوشمند
دل به ذوقِ خردهٔ مینا ببند
</poem>
